\documentclass{article}
\usepackage[T1]{fontenc}
\usepackage{lmodern}
\usepackage{graphicx}
\usepackage{amsmath,amssymb}
\usepackage[a4paper,margin=2cm]{geometry}
\usepackage[absolute,overlay]{textpos}
\setlength{\TPHorizModule}{1mm}
\setlength{\TPVertModule}{1mm}
\usepackage[indonesian]{babel}
\usepackage{hyperref}
\setlength{\emergencystretch}{2em}
% unit: Halaman 7

\begin{document}

% === Halaman 7 (python/native) ===

• Contoh pada Caesar Cipher, penyerang hanya memiliki cipherteks berikut:

                      KHOORZRUOG

   Penyerang tidak tahu kunci (pergeseran berapa huruf), hanya punya cipherteks itu 

saja. 

Namun karena Caesar cipher hanya punya 25 kemungkinan pergeseran huruf, 

penyerang bisa mencoba semua pergeseran (brute force).

Saat dicoba pergeseran 3 huruf ke kiri, diperoleh plainteks:

                 HELLOWORLD

Serangan berhasil menemukan kunci dan memecahkan cipherteks!

7

\end{document}
